% LPC 2026, Refereed track -- 6 October 2026, 10:00-10:45, Small Hall, Prague:
% "RCU Pseudo-Transactions: Bridging the gap between RCU and STM".
%
% The deck is P1 (../p1-sw-flip-latch) told to a room of kernel developers: its
% numbers (engine pin 2793224e), its framing, and its CSVs, read in place so no
% chart here can drift from the paper's. HANDOFF.md holds the outline, the
% decisions behind it and the standing rules; `make` checks the framing words.
%
% Scope, decided 2026-09-30: the single-writer API only. The CAS-based
% multi-writer engine (P2) gets no slide, not even a teaser.
%
% The dentry cache, decided 2026-10-03 once its benchmarks finished: one slide
% and one backup slide, on the bucket lock + SW txn engine only, with the
% numbers of the benchmark tree's sweep (data/, dcache-summary.py): re-taken
% 2026-10-03 at P1's engine pin, 2793224e, built without assertions.

\documentclass[aspectratio=169, 10pt, t, xcolor=table]{beamer}
% `make notes` defines \withnotes: every slide is followed by its notes page.
\ifdefined\withnotes\setbeameroption{show notes}\fi

% ---------------------------------------------------------------------------
% Fonts. IBM Plex is the design's face (HANDOFF.md section 4); where it is not
% installed, Lato + Noto Sans Mono stand in. Both cover every glyph used here,
% arrows included. A font swap can move a line break, so `make` reports any
% overfull box.
% ---------------------------------------------------------------------------
\usepackage{fontspec}
\IfFontExistsTF{IBM Plex Sans}{%
  % The 600 weight's name depends on the build: upstream's OTFs call it
  % "SemiBold", Debian's fonts-ibm-plex TTFs "SmBld".
  \IfFontExistsTF{IBM Plex Sans SemiBold}{\def\plexsb{SemiBold}}{\def\plexsb{SmBld}}%
  \setsansfont{IBM Plex Sans}[BoldFont={IBM Plex Sans \plexsb},
                              BoldItalicFont={IBM Plex Sans \plexsb\space Italic}]
  \setmonofont{IBM Plex Mono}[BoldFont={IBM Plex Mono \plexsb}]
}{%
  \setsansfont{Lato}[BoldFont={Lato Semibold},
                     BoldItalicFont={Lato Semibold Italic}]
  % One numeric scale, not Scale=MatchLowercase: fontspec computes that per
  % face (0.945 regular, 0.928 bold), which knocks the aligned comments of a
  % listing out of line wherever a bold identifier precedes them.
  \setmonofont{Noto Sans Mono}[Scale=0.945,
                               BoldFont={Noto Sans Mono Bold}]
}
% Math letters and digits in the text face: the slides' math is single letters
% ($E$, $o_1$, $s_k$), and in Computer Modern they read as a second font.
\usepackage[italic]{mathastext}

\usepackage{xspace}
% Shared macros for the urcu-txn paper series.
% Included by each paper's main.tex via % Shared macros for the urcu-txn paper series.
% Included by each paper's main.tex via % Shared macros for the urcu-txn paper series.
% Included by each paper's main.tex via \input{../common/macros}.
%
% Keep this class-agnostic: nothing here may assume `article`, so that a paper
% can be retargeted to acmart/IEEEtran by changing only its own preamble.

% ---------------------------------------------------------------------------
% Terminology.
%
% These exist so the series says the same thing the same way every time, and so
% that a terminology decision is changed in one place. The guarantee wording is
% load-bearing: this engine is NOT STM and is NOT serializable. See
% \pseudotxn below and the note in README.md.
% ---------------------------------------------------------------------------
\newcommand{\pseudotxn}{pseudo-transaction\xspace}
\newcommand{\pseudotxns}{pseudo-transactions\xspace}
\newcommand{\mcas}{MCAS\xspace}
\newcommand{\kcas}{$k$-CAS\xspace}
\newcommand{\rcu}{RCU\xspace}
\newcommand{\QSBR}{QSBR\xspace}
\newcommand{\urcu}{Userspace RCU\xspace}
\newcommand{\sw}{single-writer\xspace}
\newcommand{\mw}{multi-writer\xspace}
\newcommand{\fliplatch}{flip-latch\xspace}
\newcommand{\ryw}{read-your-own-writes\xspace}

% Engine identifiers.
\newcommand{\urcutxn}{\texttt{urcu-txn}\xspace}
\newcommand{\rcumcas}{\texttt{rcu-mcas}\xspace}
\newcommand{\rcutxnsw}{\texttt{rcu-txn-sw}\xspace}

% ---------------------------------------------------------------------------
% Code and identifiers.
% ---------------------------------------------------------------------------
% \code offers TeX no place to break, and this series is dense with identifiers
% long enough to overflow a line on their own (cds_list_for_each_entry_reverse,
% -DURCU_TXN_SW_EXCL_VALIDATE). Two of them in one sentence is enough to push
% text into the margin no matter how the surrounding prose is stretched.
%
% Since every such identifier already writes its underscores as \_, redefining
% \_ locally gives TeX a legal breakpoint after each one -- silently, with NO
% hyphen. That last part is the whole reason not to reach for hyphenat's [htt]
% option instead: a hyphen inside a C identifier is ambiguous, and the reader
% cannot tell whether it belongs to the name. For a paper whose purpose is
% enablement, an identifier someone has to guess at is a real defect; a slightly
% loose line is not.
%
% Scoped to \texttt's own group, so \_ keeps its normal meaning everywhere else.
%
% \let, not \renewcommand, and \DeclareRobustCommand, not \newcommand: \code is
% used inside \caption, which is a MOVING argument (written to the .aux and
% re-read). A fragile \code expands there and the \renewcommand form dies with
% "Argument of \caption@ydblarg has an extra }" -- a fatal error whose message
% names neither \code nor the caption it came from.
\newcommand{\codeunderscore}{\textunderscore\allowbreak}
\DeclareRobustCommand{\code}[1]{\texttt{\let\_\codeunderscore #1}}
\DeclareRobustCommand{\fn}[1]{\code{#1()}}

% A record in a transaction's write set: {slot, old, new}.
\newcommand{\record}[3]{$\{#1,\,#2 \rightarrow #3\}$}
% An edit on one slot, old to new.
\newcommand{\edit}[2]{$#1 \rightarrow #2$}

% Transaction status values.
\newcommand{\UNDECIDED}{\textsc{undecided}\xspace}
\newcommand{\SUCCEEDED}{\textsc{succeeded}\xspace}
\newcommand{\FAILED}{\textsc{failed}\xspace}
\newcommand{\ABORT}{\textsc{abort}\xspace}

% ---------------------------------------------------------------------------
% Editorial markers. Define \draftmode in the preamble to make these visible;
% they vanish in a submission build. Never ship with \draftmode set.
% ---------------------------------------------------------------------------
\newif\ifdraftmode
\draftmodefalse

% These use \color/\bfseries SWITCHES inside a group rather than
% \textcolor{}{}/\textbf{}, whose arguments are short: an editorial note
% spanning a blank line would otherwise die with "Paragraph ended before
% \@textcolor was complete". Notes routinely span paragraphs, so the switch
% form is the only one that survives real use.
\newcommand{\marker}[3]{\ifdraftmode{\color{#1}\bfseries[#2: #3]}\fi}

\newcommand{\TODO}[1]{\marker{red}{TODO}{#1}}
\newcommand{\NOTE}[1]{\marker{blue}{NOTE}{#1}}
% Claims needing a citation before submission.
\newcommand{\NEEDCITE}[1]{\marker{orange}{CITE}{#1}}
% Numbers currently sourced from source comments rather than a reproducible run.
\newcommand{\UNVERIFIED}[1]{\marker{purple}{UNVERIFIED}{#1}}
.
%
% Keep this class-agnostic: nothing here may assume `article`, so that a paper
% can be retargeted to acmart/IEEEtran by changing only its own preamble.

% ---------------------------------------------------------------------------
% Terminology.
%
% These exist so the series says the same thing the same way every time, and so
% that a terminology decision is changed in one place. The guarantee wording is
% load-bearing: this engine is NOT STM and is NOT serializable. See
% \pseudotxn below and the note in README.md.
% ---------------------------------------------------------------------------
\newcommand{\pseudotxn}{pseudo-transaction\xspace}
\newcommand{\pseudotxns}{pseudo-transactions\xspace}
\newcommand{\mcas}{MCAS\xspace}
\newcommand{\kcas}{$k$-CAS\xspace}
\newcommand{\rcu}{RCU\xspace}
\newcommand{\QSBR}{QSBR\xspace}
\newcommand{\urcu}{Userspace RCU\xspace}
\newcommand{\sw}{single-writer\xspace}
\newcommand{\mw}{multi-writer\xspace}
\newcommand{\fliplatch}{flip-latch\xspace}
\newcommand{\ryw}{read-your-own-writes\xspace}

% Engine identifiers.
\newcommand{\urcutxn}{\texttt{urcu-txn}\xspace}
\newcommand{\rcumcas}{\texttt{rcu-mcas}\xspace}
\newcommand{\rcutxnsw}{\texttt{rcu-txn-sw}\xspace}

% ---------------------------------------------------------------------------
% Code and identifiers.
% ---------------------------------------------------------------------------
% \code offers TeX no place to break, and this series is dense with identifiers
% long enough to overflow a line on their own (cds_list_for_each_entry_reverse,
% -DURCU_TXN_SW_EXCL_VALIDATE). Two of them in one sentence is enough to push
% text into the margin no matter how the surrounding prose is stretched.
%
% Since every such identifier already writes its underscores as \_, redefining
% \_ locally gives TeX a legal breakpoint after each one -- silently, with NO
% hyphen. That last part is the whole reason not to reach for hyphenat's [htt]
% option instead: a hyphen inside a C identifier is ambiguous, and the reader
% cannot tell whether it belongs to the name. For a paper whose purpose is
% enablement, an identifier someone has to guess at is a real defect; a slightly
% loose line is not.
%
% Scoped to \texttt's own group, so \_ keeps its normal meaning everywhere else.
%
% \let, not \renewcommand, and \DeclareRobustCommand, not \newcommand: \code is
% used inside \caption, which is a MOVING argument (written to the .aux and
% re-read). A fragile \code expands there and the \renewcommand form dies with
% "Argument of \caption@ydblarg has an extra }" -- a fatal error whose message
% names neither \code nor the caption it came from.
\newcommand{\codeunderscore}{\textunderscore\allowbreak}
\DeclareRobustCommand{\code}[1]{\texttt{\let\_\codeunderscore #1}}
\DeclareRobustCommand{\fn}[1]{\code{#1()}}

% A record in a transaction's write set: {slot, old, new}.
\newcommand{\record}[3]{$\{#1,\,#2 \rightarrow #3\}$}
% An edit on one slot, old to new.
\newcommand{\edit}[2]{$#1 \rightarrow #2$}

% Transaction status values.
\newcommand{\UNDECIDED}{\textsc{undecided}\xspace}
\newcommand{\SUCCEEDED}{\textsc{succeeded}\xspace}
\newcommand{\FAILED}{\textsc{failed}\xspace}
\newcommand{\ABORT}{\textsc{abort}\xspace}

% ---------------------------------------------------------------------------
% Editorial markers. Define \draftmode in the preamble to make these visible;
% they vanish in a submission build. Never ship with \draftmode set.
% ---------------------------------------------------------------------------
\newif\ifdraftmode
\draftmodefalse

% These use \color/\bfseries SWITCHES inside a group rather than
% \textcolor{}{}/\textbf{}, whose arguments are short: an editorial note
% spanning a blank line would otherwise die with "Paragraph ended before
% \@textcolor was complete". Notes routinely span paragraphs, so the switch
% form is the only one that survives real use.
\newcommand{\marker}[3]{\ifdraftmode{\color{#1}\bfseries[#2: #3]}\fi}

\newcommand{\TODO}[1]{\marker{red}{TODO}{#1}}
\newcommand{\NOTE}[1]{\marker{blue}{NOTE}{#1}}
% Claims needing a citation before submission.
\newcommand{\NEEDCITE}[1]{\marker{orange}{CITE}{#1}}
% Numbers currently sourced from source comments rather than a reproducible run.
\newcommand{\UNVERIFIED}[1]{\marker{purple}{UNVERIFIED}{#1}}
.
%
% Keep this class-agnostic: nothing here may assume `article`, so that a paper
% can be retargeted to acmart/IEEEtran by changing only its own preamble.

% ---------------------------------------------------------------------------
% Terminology.
%
% These exist so the series says the same thing the same way every time, and so
% that a terminology decision is changed in one place. The guarantee wording is
% load-bearing: this engine is NOT STM and is NOT serializable. See
% \pseudotxn below and the note in README.md.
% ---------------------------------------------------------------------------
\newcommand{\pseudotxn}{pseudo-transaction\xspace}
\newcommand{\pseudotxns}{pseudo-transactions\xspace}
\newcommand{\mcas}{MCAS\xspace}
\newcommand{\kcas}{$k$-CAS\xspace}
\newcommand{\rcu}{RCU\xspace}
\newcommand{\QSBR}{QSBR\xspace}
\newcommand{\urcu}{Userspace RCU\xspace}
\newcommand{\sw}{single-writer\xspace}
\newcommand{\mw}{multi-writer\xspace}
\newcommand{\fliplatch}{flip-latch\xspace}
\newcommand{\ryw}{read-your-own-writes\xspace}

% Engine identifiers.
\newcommand{\urcutxn}{\texttt{urcu-txn}\xspace}
\newcommand{\rcumcas}{\texttt{rcu-mcas}\xspace}
\newcommand{\rcutxnsw}{\texttt{rcu-txn-sw}\xspace}

% ---------------------------------------------------------------------------
% Code and identifiers.
% ---------------------------------------------------------------------------
% \code offers TeX no place to break, and this series is dense with identifiers
% long enough to overflow a line on their own (cds_list_for_each_entry_reverse,
% -DURCU_TXN_SW_EXCL_VALIDATE). Two of them in one sentence is enough to push
% text into the margin no matter how the surrounding prose is stretched.
%
% Since every such identifier already writes its underscores as \_, redefining
% \_ locally gives TeX a legal breakpoint after each one -- silently, with NO
% hyphen. That last part is the whole reason not to reach for hyphenat's [htt]
% option instead: a hyphen inside a C identifier is ambiguous, and the reader
% cannot tell whether it belongs to the name. For a paper whose purpose is
% enablement, an identifier someone has to guess at is a real defect; a slightly
% loose line is not.
%
% Scoped to \texttt's own group, so \_ keeps its normal meaning everywhere else.
%
% \let, not \renewcommand, and \DeclareRobustCommand, not \newcommand: \code is
% used inside \caption, which is a MOVING argument (written to the .aux and
% re-read). A fragile \code expands there and the \renewcommand form dies with
% "Argument of \caption@ydblarg has an extra }" -- a fatal error whose message
% names neither \code nor the caption it came from.
\newcommand{\codeunderscore}{\textunderscore\allowbreak}
\DeclareRobustCommand{\code}[1]{\texttt{\let\_\codeunderscore #1}}
\DeclareRobustCommand{\fn}[1]{\code{#1()}}

% A record in a transaction's write set: {slot, old, new}.
\newcommand{\record}[3]{$\{#1,\,#2 \rightarrow #3\}$}
% An edit on one slot, old to new.
\newcommand{\edit}[2]{$#1 \rightarrow #2$}

% Transaction status values.
\newcommand{\UNDECIDED}{\textsc{undecided}\xspace}
\newcommand{\SUCCEEDED}{\textsc{succeeded}\xspace}
\newcommand{\FAILED}{\textsc{failed}\xspace}
\newcommand{\ABORT}{\textsc{abort}\xspace}

% ---------------------------------------------------------------------------
% Editorial markers. Define \draftmode in the preamble to make these visible;
% they vanish in a submission build. Never ship with \draftmode set.
% ---------------------------------------------------------------------------
\newif\ifdraftmode
\draftmodefalse

% These use \color/\bfseries SWITCHES inside a group rather than
% \textcolor{}{}/\textbf{}, whose arguments are short: an editorial note
% spanning a blank line would otherwise die with "Paragraph ended before
% \@textcolor was complete". Notes routinely span paragraphs, so the switch
% form is the only one that survives real use.
\newcommand{\marker}[3]{\ifdraftmode{\color{#1}\bfseries[#2: #3]}\fi}

\newcommand{\TODO}[1]{\marker{red}{TODO}{#1}}
\newcommand{\NOTE}[1]{\marker{blue}{NOTE}{#1}}
% Claims needing a citation before submission.
\newcommand{\NEEDCITE}[1]{\marker{orange}{CITE}{#1}}
% Numbers currently sourced from source comments rather than a reproducible run.
\newcommand{\UNVERIFIED}[1]{\marker{purple}{UNVERIFIED}{#1}}

\usepackage{booktabs}
\usepackage{tikz}
\usetikzlibrary{positioning, arrows.meta, calc, fit, backgrounds,
                decorations.pathreplacing}
\usepackage{pgfplots}
\pgfplotsset{compat=1.18}
% Numbers in the deck's text face, not in math-mode Computer Modern: pgf wraps
% every printed number in $...$ unless told it is already in math mode, and
% pgfplots' default tick labels add their own $...$ on top.
\pgfkeys{/pgf/number format/assume math mode=true}
\pgfplotsset{
  xticklabel={\pgfmathprintnumber{\tick}},
  yticklabel={\pgfmathprintnumber{\tick}},
}
\usepgfplotslibrary{groupplots}
\usepackage{pgfplotstable}
\usepackage{listings}
\usepackage{tcolorbox}
\tcbuselibrary{listings, skins}

% P1's plot data, read in place.
\newcommand{\ponedata}{../p1-sw-flip-latch/data}

% The dentry-cache slide's ratios (dcache-summary.py writes them from the
% benchmark tree): \dcrange{key} is a row's range as the chart prints it,
% \dcval{key}{lo|med|hi} one end of it, \dcacheSweep the sweep's provenance id.
% Written by lpc2026-slides/dcache-summary.py; do not edit.
% efficios-trie-benchmark 3a79ed3, sweep c03704139064-c21f5a38:
% bucket lock + SW txn / seqlock baseline, as that tree's
% scripts/plot_dcache_bucketlock_summary.py computes it.
\def\dcacheSweep{c03704139064-c21f5a38}
\def\dcacheBench{3a79ed3}
\expandafter\def\csname dcv/lookup-low/lo\endcsname{1.03}
\expandafter\def\csname dcv/lookup-low/med\endcsname{1.04}
\expandafter\def\csname dcv/lookup-low/hi\endcsname{1.05}
\expandafter\def\csname dcv/lookup-rest/lo\endcsname{0.98}
\expandafter\def\csname dcv/lookup-rest/med\endcsname{1.01}
\expandafter\def\csname dcv/lookup-rest/hi\endcsname{1.05}
\expandafter\def\csname dcv/churn-alloc/lo\endcsname{0.99}
\expandafter\def\csname dcv/churn-alloc/med\endcsname{1.00}
\expandafter\def\csname dcv/churn-alloc/hi\endcsname{1.04}
\expandafter\def\csname dcv/lookup-high/lo\endcsname{1.20}
\expandafter\def\csname dcv/lookup-high/med\endcsname{1.51}
\expandafter\def\csname dcv/lookup-high/hi\endcsname{1.83}
\expandafter\def\csname dcv/dpath/lo\endcsname{1.41}
\expandafter\def\csname dcv/dpath/med\endcsname{1.52}
\expandafter\def\csname dcv/dpath/hi\endcsname{7.16}
\expandafter\def\csname dcv/readdir/lo\endcsname{1.11}
\expandafter\def\csname dcv/readdir/med\endcsname{1.58}
\expandafter\def\csname dcv/readdir/hi\endcsname{3.67}
\expandafter\def\csname dcv/churn-inplace/lo\endcsname{1.14}
\expandafter\def\csname dcv/churn-inplace/med\endcsname{1.17}
\expandafter\def\csname dcv/churn-inplace/hi\endcsname{1.23}
\expandafter\def\csname dcv/hit/lo\endcsname{0.84}
\expandafter\def\csname dcv/hit/med\endcsname{0.87}
\expandafter\def\csname dcv/hit/hi\endcsname{0.93}
\expandafter\def\csname dcv/hit-rest/lo\endcsname{0.90}
\expandafter\def\csname dcv/hit-rest/med\endcsname{0.91}
\expandafter\def\csname dcv/hit-rest/hi\endcsname{0.95}
\expandafter\def\csname dcv/cap-leaf/lo\endcsname{20.77}
\expandafter\def\csname dcv/cap-leaf/med\endcsname{20.97}
\expandafter\def\csname dcv/cap-leaf/hi\endcsname{21.48}
\expandafter\def\csname dcv/cap-dir/lo\endcsname{4.06}
\expandafter\def\csname dcv/cap-dir/med\endcsname{4.19}
\expandafter\def\csname dcv/cap-dir/hi\endcsname{4.48}

\newcommand{\dcval}[2]{\ifcsname dcv/#1/#2\endcsname
  \csname dcv/#1/#2\endcsname\else
  \errmessage{no dcache ratio `#1/#2' in data/dcache-summary-macros.tex}\fi}
\newcommand{\dcrange}[1]{\dcval{#1}{lo}--\dcval{#1}{hi}×}

% \lastof{table}{column}{\macro}: \macro <- the column's value in the table's
% last row. Every P1 CSV ends at 192 workers (191 where one core runs the
% writer), the point the talk quotes.
\newcommand{\lastof}[3]{%
  \pgfplotstablegetrowsof{#1}%
  \pgfmathtruncatemacro{\lastrow}{\pgfplotsretval-1}%
  \pgfplotstablegetelem{\lastrow}{#2}\of#1\let#3\pgfplotsretval}

% ---------------------------------------------------------------------------
% Palette (HANDOFF.md section 4; contrasts checked there). Orange is plain RCU
% and "old", blue is the facility and "new"; the whole deck keeps to that.
% ---------------------------------------------------------------------------
\definecolor{paper}{HTML}{F5F3EE}
\definecolor{paperalt}{HTML}{E9E5DC}
\definecolor{dark}{HTML}{17191E}
\definecolor{darkmuted}{HTML}{B8BDC7}
\definecolor{ink}{HTML}{1A1D23}
\definecolor{muted}{HTML}{4F5561}
\definecolor{accent}{HTML}{2456A6}
\definecolor{old}{HTML}{B04A0C}
\definecolor{codebg}{HTML}{1F232B}
\definecolor{codefg}{HTML}{E8E4DA}
\definecolor{codecomment}{HTML}{9AA3B2}
\definecolor{codekw}{HTML}{F2A65A}
\definecolor{codeblue}{HTML}{8DB6F0}
% writer-scaling series, CVD-checked as a set (and told apart by dash too)
\definecolor{stripes}{HTML}{6A8BC6}
\definecolor{onelock}{HTML}{707682}

% ---------------------------------------------------------------------------
% Beamer look: quiet, code-first. A small section kicker above each title
% stands in for divider slides, which a 45-minute slot cannot spare.
% ---------------------------------------------------------------------------
\setbeamersize{text margin left=9mm, text margin right=9mm}
\setbeamertemplate{navigation symbols}{}
\setbeamercolor{background canvas}{bg=paper}
\setbeamercolor{normal text}{fg=ink, bg=paper}
\setbeamercolor{structure}{fg=accent}
\setbeamercolor{alerted text}{fg=accent}
\setbeamercolor{frametitle}{fg=ink}
\setbeamercolor{kicker}{fg=muted}
\setbeamercolor{footline}{fg=muted}
\setbeamerfont{frametitle}{size=\LARGE, series=\bfseries}
\setbeamerfont{kicker}{size=\scriptsize, series=\bfseries}
\setbeamerfont{footline}{size=\tiny}

% \talksection sets both beamer's section and the kicker. \gdef, and a plain
% \def for the empty case: the template tests it against \@empty with \ifx,
% which a \newcommand-defined (\long) macro would never equal.
\def\kickertext{}
\newcommand{\talksection}[1]{\section{#1}\gdef\kickertext{#1}}

\makeatletter
\setbeamertemplate{frametitle}{%
  \vskip 4mm
  \ifx\kickertext\@empty\else
    {\usebeamerfont{kicker}\usebeamercolor[fg]{kicker}%
     \addfontfeature{LetterSpace=6}\MakeUppercase{\kickertext}\par}%
    \vskip 1mm
  \fi
  {\usebeamerfont{frametitle}\usebeamercolor[fg]{frametitle}%
   \insertframetitle\par}%
  \vskip 3mm}
\makeatother

% The EfficiOS wordmark, for the footline. efficios-logo.pdf is vector art on
% a letter page, and the viewport boxes it: left and right at the ink's edges,
% top at the top of the E, bottom at the letters' baseline (y = 424.1 bp; the
% ff descends to 405.6). It is not clipped, so the descenders hang under the
% box as a glyph's would; height= is therefore the height of the E.
%
% Where and how big is set by the slides, not by taste: body text runs to
% within 2.5 mm of the footline's baseline at the left margin (slides 22, 25,
% 29), and to 3.2 mm at the right. So the logo sits at the right, by the
% frame number, 3 mm tall, lowered 0.85 mm to centre on the footline's
% capitals. \smash: the footline stays as tall as its text, so the logo moves
% no slide's layout. Black on transparent: it needs a light background (the
% dark cover is a plain frame and has no footline).
\newcommand{\footlogo}{\smash{\raisebox{-0.85mm}{%
  \includegraphics[viewport=149.4 424.1 457.3 501.2, height=3mm]%
    {efficios-logo.pdf}}}}

% The frame number is set in a box as wide as two digits: the logo beside it
% must not shift between slides 9 and 10.
\setbeamertemplate{footline}{%
  \usebeamerfont{footline}\usebeamercolor[fg]{footline}%
  \hbox to \paperwidth{\hskip 9mm RCU Pseudo-Transactions\enspace·\enspace
    LPC 2026\hfill\footlogo\hskip 2.5mm
    \hbox to 2.6mm{\hfil\insertframenumber}\hskip 9mm}%
  \vskip 3.5mm}

\setbeamertemplate{itemize item}{\textcolor{accent}{–}}
\setbeamertemplate{itemize subitem}{\textcolor{muted}{–}}
\setbeamertemplate{enumerate item}{\textcolor{accent}{\insertenumlabel.}}
\setlength{\leftmargini}{1.1em}

% ---------------------------------------------------------------------------
% Code panels. The facility's identifiers are blue, C keywords orange, RCU
% primitives bold: the eye should find the API at a glance.
% ---------------------------------------------------------------------------
\lstdefinestyle{slide}{
  language=C,
  basicstyle=\ttfamily\footnotesize\color{codefg},
  keywordstyle=\color{codekw},
  keywordstyle=[2]\color{codeblue},
  keywordstyle=[3]\bfseries\color{codefg},
  commentstyle=\color{codecomment},
  stringstyle=\color{codefg},
  identifierstyle=\color{codefg},
  columns=fullflexible, keepspaces=true, showstringspaces=false,
  tabsize=4,
  morekeywords={uintptr_t},
  morekeywords=[2]{urcu_txn_sw_init, urcu_txn_sw_reserve, urcu_txn_sw_load,
    urcu_txn_sw_record, urcu_txn_sw_record_chain, urcu_txn_sw_commit,
    urcu_txn_sw_resolve, urcu_txn_sw_is_proxy, urcu_txn_sw_proxy_get,
    urcu_txn_sw_untag, urcu_txn_sw_txn, urcu_txn_sw_list_pending,
    urcu_txn_sw_list_node, urcu_txn_sw_list_add_after_prepare,
    urcu_txn_sw_list_for_each_entry_reverse_rcu, urcu_txn_sw_list_del_rcu,
    urcu_txn_sw_hlist_add_head_prepare, URCU_TXN_SW_LIST_PROXY_TAG,
    URCU_TXN_STATUS_OK, URCU_TXN_STATUS_MEMORY_ERROR},
  morekeywords=[3]{rcu_read_lock, rcu_read_unlock, rcu_dereference,
    rcu_assign_pointer, call_rcu, WRITE_ONCE, list_next_rcu, uatomic_load,
    CMM_ACQUIRE},
}
\tcbset{codebox/.style={listing only, colback=codebg, colframe=codebg,
  boxrule=0pt, arc=2pt, left=2.5mm, right=2.5mm, top=1.2mm, bottom=1.2mm,
  before skip=1mm, after skip=1mm}}
\newtcblisting{codeblock}[1][]{codebox, listing options={style=slide}, #1}
% For the few listings whose lines are real kernel or header code and too long
% to re-wrap without misquoting them.
\newtcblisting{codeblocksmall}[1][]{codebox,
  listing options={style=slide, basicstyle=\ttfamily\scriptsize\color{codefg}}, #1}

% ---------------------------------------------------------------------------
% Text blocks.
% ---------------------------------------------------------------------------
\newcommand{\lead}[1]{{\large #1\par}}
\newcommand{\punch}[1]{{\large\bfseries #1\par}}
\newcommand{\aside}[1]{{\footnotesize\color{muted}#1\par}}
% A step number in a table (defined here: a \def with a parameter inside a
% non-fragile frame would need its # doubled).
\newcommand{\stepnum}[1]{\textcolor{accent}{\textbf{#1}}}
\newtcolorbox{callout}[1][]{enhanced, colback=paperalt, frame hidden, arc=2pt,
  borderline west={1.6pt}{0pt}{accent}, left=3mm, right=3mm, top=1.5mm,
  bottom=1.5mm, fontupper=\small, before skip=1mm, after skip=1mm, #1}

% Placeholders: dashed orange, impossible to miss on screen, and listed by
% `make` so they are impossible to miss at build time too.
\newtcolorbox{placeholderbox}{enhanced, colback=old!5, frame hidden, arc=2pt,
  borderline={0.9pt}{0pt}{old, dashed}, left=3mm, right=3mm, top=2mm,
  bottom=2mm, fontupper=\small\color{old}}
\newcommand{\placeholder}[1]{%
  \begin{placeholderbox}\textbf{PLACEHOLDER}\enspace #1\end{placeholderbox}}
\newcommand{\ph}[1]{\tikz[baseline=(p.base)]\node[draw=old, dashed, text=old,
  rounded corners=1pt, inner sep=1.5pt, font=\footnotesize] (p) {#1};}

% A centred figure. A command rather than a bare {\centering ...} group: beamer
% reads a brace group right after the frame title as the frame's SUBTITLE, and
% a figure opening a frame vanished that way (this template shows no subtitle).
\newcommand{\centerfig}[1]{\par{\centering\input{#1}\par}}

% A dark frame: cover and close.
\newenvironment{darkslide}{%
  \setbeamercolor{background canvas}{bg=dark}%
  \setbeamercolor{normal text}{fg=paper, bg=dark}%
  \setbeamercolor{frametitle}{fg=paper}%
  \setbeamercolor{kicker}{fg=darkmuted}%
  \setbeamercolor{footline}{fg=darkmuted}}{}

\title{RCU Pseudo-Transactions}
\subtitle{Bridging the gap between RCU and STM}
\author{Mathieu Desnoyers}
\institute{EfficiOS Inc.}
\date{Linux Plumbers Conference 2026}

\begin{document}

% ===========================================================================
% Cover
% ===========================================================================
\begin{darkslide}
\begin{frame}[plain]
  \color{paper}%
  \vspace*{16mm}
  {\color{codeblue}\rule{16mm}{1.1mm}\par}
  \vspace{5mm}
  {\fontsize{26}{31}\selectfont\bfseries RCU Pseudo-Transactions\par}
  \vspace{2.5mm}
  {\Large\color{darkmuted}Bridging the gap between RCU and STM\par}
  \vfill
  {\large\bfseries Mathieu Desnoyers\enspace
   {\mdseries\color{darkmuted}·\enspace EfficiOS}\par}
  \vspace{1.5mm}
  {\small\color{darkmuted}Linux Plumbers Conference 2026\enspace·\enspace
   Prague\enspace·\enspace 6 October 2026\par}
  \vspace*{11mm}
  \note{%
  \begin{itemize}
  \item 45 minutes including questions; about 33 minutes of slides, 1.3 min
        per slide. Slides marked [opt] in these notes go first if short.
  \item The title's bridge is positioning, not a claim of STM semantics:
        say ``between'', never ``is''.
  \end{itemize}}
\end{frame}
\end{darkslide}

% ===========================================================================
% A. The problem
% ===========================================================================
\talksection{The problem}

\begin{frame}{RCU's contract}
  \lead{Readers pay nothing. In exchange, a writer gets \textbf{exactly one}
  atomic operation: \textbf{a single pointer store}.}
  \vspace{7mm}
  \centerfig{fig-contract}
  \vspace{7mm}
  \lead{A pseudo-transaction does not add a fourth phase.\\
  It \textbf{\color{accent}widens the second one}: N slots, still one store.}
  \note{%
  \begin{itemize}
  \item Build owes readers nothing: no ordering, no barrier; a failed build
        frees its work at once. Reclaim is merely deferred.
  \item Publish is the constraint: a store is the only thing commodity
        hardware makes atomic to a reader that took no lock and paid no barrier.
  \item The rest of the talk is about the second half of the first sentence.
  \end{itemize}}
\end{frame}

\begin{frame}[fragile]{The kernel's list already needs two stores}
\begin{codeblocksmall}
static inline void __list_del(struct list_head *prev,
                              struct list_head *next)
{
    next->prev = prev;              /* N->prev: E -> P   plain store  */
    WRITE_ONCE(prev->next, next);   /* P->next: E -> N   marked store */
}

/* __list_add_rcu(): */
new->next = next;                              /* private */
new->prev = prev;                              /* private */
rcu_assign_pointer(list_next_rcu(prev), new);  /* publishes new */
next->prev = new;                              /* plain store, AFTER publish */
\end{codeblocksmall}
\vspace{2mm}
\begin{itemize}
\item \code{list\_del\_rcu()} poisons \code{prev}: rculist ships forward
      iterators only --- and so does Userspace RCU's \code{cds\_list\_*\_rcu}.
\item A considered trade, not an oversight: the poison buys use-after-delete
      debugging.
\end{itemize}
\note{%
\begin{itemize}
\item Between the last two lines of \_\_list\_add\_rcu(): prev->next == new
      while next->prev == prev. No weak-memory argument needed: it is program
      order, so the window is real on every architecture.
\item Say the trade out loud, and mean it: it earns goodwill, and it is true.
      The kernel judged forward-walk debugging worth more than a reverse walk;
      Userspace RCU mirrored the kernel's interface.
\end{itemize}}
\end{frame}

\begin{frame}{Publication atomicity, not ordering}
\begin{columns}[T, onlytextwidth]
\column{0.45\textwidth}
  {\centering% Removing E from P <-> E <-> N with BOTH edges properly published (P1
% sec:onestore, "Ordering is not what is missing"), drawn at the instant
% between the two release stores: the forward chain has dropped E, the backward
% chain still has it. The ghost's own pointers are what lets a reader standing
% on E escape -- and what keeps the backward view trailing for a grace period.
\begin{tikzpicture}[
  lnode/.style = {draw=ink, fill=paperalt, rounded corners=2pt,
                  minimum width=11mm, minimum height=7.5mm, font=\ttfamily},
  ghost/.style = {lnode, dashed, draw=muted, fill=paper, text=muted},
  edge/.style  = {-{Latex[length=2.2mm]}, line width=0.9pt},
  own/.style   = {-{Latex[length=1.8mm]}, densely dashed, muted, line width=0.6pt},
  lab/.style   = {font=\footnotesize\ttfamily},
]
\node[lnode] (P) at (0,0)   {P};
\node[ghost] (E) at (2.5,0) {E};
\node[lnode] (N) at (5.0,0) {N};
% the ghost keeps its own next and prev
\draw[own] ([yshift=1mm]E.east) -- ([yshift=1mm]N.west);
\draw[own] ([yshift=-1mm]E.west) -- ([yshift=-1mm]P.east);
% first store: done
\draw[edge, ink] (P.north) to[out=40, in=140]
  node[above, lab] {P->next == N} (N.north);
% second store: not yet
\draw[edge, old, line width=1.3pt] (N.south) to[out=-140, in=-40]
  node[below, lab, text=old] {N->prev == E} (E.south);
\end{tikzpicture}
\par}
  \vspace{3mm}
  \punch{No barrier fixes a window whose width is a grace period.}
\column{0.52\textwidth}
  \small
  \begin{itemize}
  \item Even properly published, two edges are \textbf{two release stores}.
        Between them the forward chain has dropped E, the backward chain has
        not: \textbf{two different lists}, neither torn.
  \item E is a ghost until its grace period completes: the backward view trails by
        \textbf{up to a grace period}.
  \end{itemize}
  \vspace{2mm}
\begin{callout}[fontupper=\footnotesize]
\textbf{The kernel's narrow exception} (Dec 2024, Brauner; suggested and
reviewed by McKenney): \code{list\_bidir\_del\_rcu()} +
\code{list\_bidir\_prev\_rcu()} let a reader step back out of a ghost ---
useful, and not a coherent reverse walk. No mixing with
\code{list\_del\_rcu()} on one list; its sole caller takes one hop.
\end{callout}
\end{columns}
\note{%
\begin{itemize}
\item The tempting reading: publish prev with rcu\_assign\_pointer(), read it
      with rcu\_dereference(), done. It is not: two release stores are two
      publication instants.
\item What is missing is atomicity across two slots. That is why the fix has
      to be a mechanism rather than a barrier.
\item The sole in-tree caller as of next-20260226:
      \_\_ns\_tree\_adjoined\_rcu() in kernel/nstree.c, one hop, prev \emph{or}
      next, never comparing the two directions.
\end{itemize}}
\end{frame}

\begin{frame}{Composition: two structures, one object}
  A cache entry $E$ must be \textbf{findable} (hash table) and
  \textbf{ordered} (LRU list).
  \vspace{2mm}
  \centerfig{fig-skew}
  \vspace{1mm}
  A writer-side lock does not help: \textbf{the observer is a reader, and
  takes no lock.}
  \label{fr:skew}
\note{%
  \begin{itemize}
  \item Publishing into the list first gives the mirror image: ordered, but
        a lookup reports it absent.
  \item The window is not a few instructions: the writer can be preempted
        inside it. Callers grow a validity flag or a lock, and the reason RCU
        was chosen erodes.
  \end{itemize}}
\end{frame}

\begin{frame}{Today's remedies each pay somewhere}
\begingroup\small\renewcommand{\arraystretch}{1.1}\arrayrulecolor{muted!60}
% A light rule between rows: several cells wrap, and without one the rows
% read as one block.
\def\rowsep{\arrayrulecolor{muted!25}\specialrule{0.4pt}{0.9mm}{0.9mm}%
  \arrayrulecolor{muted!60}}
\begin{tabular}{@{}>{\raggedright}p{38mm}>{\raggedright}p{22mm}>{\raggedright\arraybackslash}p{72mm}@{}}
\toprule
\textbf{Remedy} & \textbf{Who pays} & \textbf{How} \\
\midrule
Copy the region & the writer
  & allocation, footprint and reclaim on every update; a pinned node cannot
    be copied at all \\
\rowsep
Drain with a grace period & the writer
  & waits out every reader in the system \\
\rowsep
Steer or validate the reader & \textbf{every reader}
  & seqcount retries (dcache: \code{d\_seq}, \mbox{\code{rename\_lock}}),
    per-element locks: paid on \textbf{every traversal} \\
\rowsep
Decline the operation & the user & no RCU reverse iterator \\
\bottomrule
\end{tabular}\par\endgroup
\vspace{4mm}
\lead{What is missing is \textbf{publication atomicity} across slots, and RCU
has no way to ask for it.}
\note{%
\begin{itemize}
\item Pinned: an address held somewhere the writer cannot reach --- an
      application handle, another subsystem's cache. A swap is complete only
      if every reference is redirected; one reference the writer does not own
      is enough to rule copying out.
\item The seqcount detects after the fact and retries; the per-element lock
      prevents the interleaving. Both are paid on the read side whether or not
      a writer ever appears.
\end{itemize}}
\end{frame}

% ===========================================================================
% B. The model: the definition (what it guarantees follows the API)
% ===========================================================================
\talksection{The model}

\begin{frame}{A pseudo-transaction}
\begin{callout}
A \textbf{frozen} set of \{\textit{slot}, \textit{old} → \textit{new}\}
records. A sequence of \textbf{install} stores makes them visible, slot by
slot, each still resolving \textit{old}; \textbf{one atomic commit step} then
publishes the transition to \textit{new}, for all of them at once.
\end{callout}
\vspace{0.5mm}
\begingroup\small\renewcommand{\arraystretch}{1.0}\arrayrulecolor{muted!60}
% The first three rows are a sequence and are numbered as one; the rows
% under the rule are properties of the result, not steps.
\begin{tabular}{@{}l@{\hspace{2.5mm}}ll@{}}
\toprule
& \textbf{Journaling filesystem} & \textbf{Pseudo-transaction} \\
\midrule
\stepnum{1} & \textbf{journal transaction closed} & \textbf{record set frozen} \\
\stepnum{2} & \textbf{journal blocks written}
         & \textbf{records installed}: slots parked, still \emph{old} \\
\stepnum{3} & \textbf{commit record written}
         & \textbf{one commit step}: \emph{old} → \emph{new}, for all \\
\midrule
& before commit: old state      & before commit: every record resolves \emph{old} \\
& after commit: new state       & after commit: every record resolves \emph{new} \\
& never a partial batch         & no slot, and no instant, is ever torn \\
& readers are not transactional & readers are plain RCU readers \\
& atomicity, not isolation      & atomicity, not opacity \\
\bottomrule
\end{tabular}\par\endgroup
\note{%
\begin{itemize}
\item Freezing settles \emph{what} takes effect, and comes first; the single
      step settles \emph{when}, and comes last.
\item Install is not a partial commit: readers can reach the records, but
      every parked slot still resolves to its old value. Visibility and the
      old-to-new transition are two different events; only the second is
      the linearization point.
\item A journaling filesystem calls this a transaction and has never meant
      isolation by it. This table is what stops the room filing the model
      under ``weak STM'' in the first thirty seconds.
\end{itemize}}
\end{frame}

% ===========================================================================
% C. The mechanism
% ===========================================================================
\talksection{The mechanism}

\begin{frame}{The flip-latch: one level of indirection}
\centerfig{fig-fliplatch}
\lead{Every proxy reads one selector: \textbf{one release store flips the
set}.}
\note{%
\begin{itemize}
\item A slot pointing at a proxy is \emph{parked}. Installing a proxy and
      parking a slot are the same event, from two sides.
\item Two states are all this facility needs: under exclusion a commit
      cannot fail for contention, and one that fails for memory does so
      before it parks anything.
\item Lineage, to say if asked and not to claim: a descriptor of records plus
      a status word that decides the outcome is practical MCAS (Harris,
      Fraser, Pratt, DISC 2002; Fraser's thesis, 2004). Install, commit and
      settle are its acquire, decision and release phases, and it too marks a
      slot by reserved low bits. P1 sec.\ 11 claims none of that. What
      exclusion buys is that no step here is a compare-and-swap, and readers
      never help.
\end{itemize}}
\end{frame}

\begin{frame}{Lifecycle: only the commit is visible}
\begingroup\renewcommand{\arraystretch}{1.35}\arrayrulecolor{muted!60}
\begin{tabular}{@{}llll@{}}
\toprule
\textbf{Step} & \textbf{The slot holds} & \textbf{A reader resolves}
  & \textbf{Linearization point?} \\
\midrule
(before) & $o$, directly & $o$ & --- \\
Build    & $o$, directly & $o$ & no: new nodes unreachable \\
Install  & tagged proxy, \code{sel} = 0 & $o$ & no: resolved value unchanged \\
\rowcolor{accent!10}
\textbf{Commit} & tagged proxy, \code{sel} = 1 & $\boldsymbol{n}$
  & \textbf{yes: one release store} \\
Settle   & $n$, directly & $n$ & no: resolved value unchanged \\
(after)  & $n$, directly & $n$ & --- \\
\bottomrule
\end{tabular}\par\endgroup
\vspace{6mm}
\begin{itemize}
\item Install and settle both rewrite the slot, and change nothing a reader
      resolves.
\item Settle \emph{is} the uninstall. Then \code{call\_rcu()} the txn
      descriptor.
\end{itemize}
\note{%
\begin{itemize}
\item The reader's column reads o, o, o, n, n, n: one transition, however many
      slots and however long the writer takes.
\item The wrinkle is the contract in disguise: across the install window a
      proxy resolves to the old value recorded at install. Under exclusion
      that is the true old. Without it, a second writer's store is silently
      overwritten by settle.
\item ``Why settle at all?'' Guerraoui, Kogan, Marathe, Zablotchi (DISC 2020)
      do not, on their critical path: the descriptor stays in the slot until a
      later operation or reclamation, which saves them $k$ CASes and costs
      their readers an indirection after every commit. Here settle is a plain
      store, so there is little to save, and settling at once is what keeps
      the steady state free of proxies.
\end{itemize}}
\end{frame}

\begin{frame}[fragile]{The read side is one branch}
\begin{codeblocksmall}
static inline void *urcu_txn_sw_resolve(void *v, uintptr_t tag)
{
    if (caa_likely(!urcu_txn_sw_is_proxy(v, tag)))  /* (v & tag) != tag */
        return v;
    return urcu_txn_sw_proxy_get(urcu_txn_sw_untag(v, tag));
}

/* urcu_txn_sw_proxy_get(): */
    proxy->ptr[uatomic_load(&proxy->group->selector, CMM_ACQUIRE)]
\end{codeblocksmall}
\vspace{3mm}
\begin{columns}[T, onlytextwidth]
\column{0.48\textwidth}
  {\bfseries\color{accent}Not parked}\enspace
  {\small\color{muted}almost always\par}
  \vspace{1mm}
  {\small A test on a register already loaded, and a predicted branch.
  \textbf{No load.}\par}
\column{0.48\textwidth}
  {\bfseries Parked}\enspace
  {\small\color{muted}install → settle, that slot only\par}
  \vspace{1mm}
  {\small proxy → group → selector → \code{ptr[sel]}: indexed, a data
  dependency, not a branch.\par}
\end{columns}
\vspace{5mm}
\aside{The untag \emph{subtracts} the tag, so the compiler folds it into the
next load's displacement.}
\note{%
\begin{itemize}
\item Only the selector is loaded with acquire, and only when parked; the
      proxy's immutable fields are reached through the dependency.
\item Subtraction is exact: the test has proven every tag bit set, and the
      record's own address has them clear.
\end{itemize}}
\end{frame}

\begin{frame}{Tags: the marker rides the pointer}
\centerfig{fig-tagword}
\vspace{5mm}
\begin{itemize}\small
\item \textbf{Per-record tag.} One commit may span structures with different
      encodings --- bit 0 here, a reserved nibble there. Cross-structure
      atomicity forces it.
\item Records are 16-byte aligned: 4 low bits to choose from. A record stores
      the tag, not a tagged pointer, so the descriptor can grow by moving.
\item \textbf{Contract}: no legitimate slot value may carry the tag pattern.
      Not checkable for free; the exclusion validator also trips on an alias
      when it parks the slot.
\end{itemize}
\note{%
\begin{itemize}
\item The tag's meaning belongs to the slot's structure: a reader never reads
      the tag from the value, it knows statically which structure it walks.
\item The record keeps its copy for the write side only: the generic commit
      parks each proxy with the pattern that slot's readers test.
\item Violating the contract is a wild pointer, not a torn read: the
      contracts underwrite memory safety, not merely atomicity.
\end{itemize}}
\end{frame}

\begin{frame}{Exclusion is per slot}
\lead{The facility owns no lock. Exclusion is the embedder's, and it is
\textbf{per slot}, not per structure.}
\vspace{4mm}
\begin{itemize}
\item \textbf{No install CAS, no conflict detection, no abort}: that is why it
      is cheap.
\item Two writers on one slot \textbf{silently corrupt} it.
      \code{-DURCU\_TXN\_SW\_EXCL\_VALIDATE} aborts the process on a violation,
      and costs nothing when off.
\item Fine-grained locks --- per bucket, per node --- are the ordinary way to
      meet it: \textbf{writers on disjoint slots commit concurrently}.
\item A commit fails only for want of memory, before anything is parked.
      \code{urcu\_txn\_sw\_reserve()} up front rules that out.
\end{itemize}
\label{fr:exclusion}
\note{%
\begin{itemize}
\item The validator claims the handle (catching a transaction handed between
      threads mid-flight) and the slot (where two racing writers collide).
      Off by default; the handle does not even carry the owner field then.
\item ``We already offer fine-grained locking with the single-writer API'':
      slide~\ref{fr:writerscale} is the evidence.
\end{itemize}}
\end{frame}

% ===========================================================================
% D. The API
% ===========================================================================
\talksection{The API}

\begin{frame}[fragile]{The API on one slide}
\begin{codeblock}
struct urcu_txn_sw_txn txn;                    /* on-stack handle */
urcu_txn_sw_init(&txn);
urcu_txn_sw_reserve(&txn, 3);                  /* rules out OOM */
old = urcu_txn_sw_load(&txn, &a->next, TAG);   /* sees own writes */
urcu_txn_sw_record_chain(&txn, &a->next, old, x, TAG);
urcu_txn_sw_record(&txn, &b->prev, y, x, TAG); /* blind append */
if (urcu_txn_sw_commit(&txn) != URCU_TXN_STATUS_OK)
    /* URCU_TXN_STATUS_MEMORY_ERROR: nothing was published */;
/* reader, under rcu_read_lock(): */
v = urcu_txn_sw_resolve(rcu_dereference(*slot), TAG);
\end{codeblock}
\vspace{1mm}
\begin{itemize}\small
\item \code{record()} appends blindly: slots pairwise distinct.
      \code{record\_chain()} fuses onto this transaction's record for the
      slot: committed \emph{old} kept, \emph{new} advanced.
\item One record commits as a plain release store; more: park, flip,
      settle, \mbox{\code{call\_rcu()}}. OOM is sticky; never an abort.
\end{itemize}
\note{%
\begin{itemize}
\item Names verified at 2793224e, include/urcu/rcu-txn-sw.h.
\item record() is a blind append over pairwise-distinct slots;
      record\_chain() fuses a second edit of a slot onto its record, else
      appends. The find is a linear scan of the private record array, with a
      Bloom filter armed lazily for long write sets. On a handle declared
      disjoint, the find is skipped.
\item The status enum is shared vocabulary; this engine returns OK or
      MEMORY\_ERROR only.
\end{itemize}}
\end{frame}

% ===========================================================================
% B, continued. The guarantees: what the model provides, and where it stands
%
% After the mechanism and the API since 2026-10-04: for this room the
% technical achievement comes first, and its place in the state of the art
% second. These three slides lean on the selector, the one branch and the
% exclusion contract, all of which the room has now seen. Their own kicker,
% so that "The model" does not come back after "The API".
% ===========================================================================
\talksection{The guarantees}

\begin{frame}{What it provides, and what it does not}
\begin{columns}[T, onlytextwidth]
\column{0.48\textwidth}
  {\large\bfseries\color{accent}Provided\par}
  \vspace{1mm}
  \begin{itemize}\small
  \item \textbf{Write-set atomicity}: every record resolves \emph{old} before
        the commit, \emph{new} after.
  \item \textbf{No tearing}: no slot, and no instant, is ever torn.
  \item \textbf{The commit linearizes at one point} --- a claim about the
        write, and only the write.
  \item \textbf{Cross-structure atomicity}. The unit is one
        pseudo-transaction: two commits are two linearization points.
  \end{itemize}
\column{0.48\textwidth}
  {\large\bfseries\color{old}Not provided\par}
  \vspace{1mm}
  \setbeamertemplate{itemize item}{\textcolor{old}{–}}
  \begin{itemize}\small
  \item \textbf{A reader is not a transaction of any kind}: a sequence of
        single-slot reads. A traversal may straddle a commit.
  \item \textbf{No opacity, no multi-read snapshot.}
  \item \textbf{No read validation}: under exclusion there is nothing to
        validate against.
  \end{itemize}
\end{columns}
\label{fr:provides}
\note{%
\begin{itemize}
\item A single-slot read is itself linearizable, at its own load. A traversal
      of several slots is a sequence of loads, not one operation.
\item Serializability is claimed nowhere. Linearizable is scoped to the
      commit, never to a read-then-commit bracket.
\item An embedder that needs a multi-read snapshot layers its own versioning
      on top (seqcount, validating transaction).
\end{itemize}}
\end{frame}

\begin{frame}{Never new-then-old}
\lead{The selector is written once, 0 → 1, and never back. A reader that
resolves the slots of one commit \textbf{along a dependency chain} sees}
\vspace{4mm}
{\centering
{\ttfamily\Large\textcolor{old}{o o o}\ \textcolor{accent}{n n n}}%
\enspace{\small\color{muted}allowed}\qquad\qquad
{\ttfamily\Large\textcolor{accent}{n}\ \textcolor{old}{o}}%
\enspace{\small\color{old}never}\par}
\vspace{6mm}
The chain is an address dependency --- every ordinary walk down a list or a
tree --- or an acquire.

\vspace{8mm}
\punch{New-then-old is what drives readers to seqcount retries. Here it
cannot happen, by construction: no seqcount, no retry.}
\note{%
\begin{itemize}
\item A reader may lag a commit; it never regresses across one.
\item Not a seqcount, and not a seqcount's guarantee: a seqcount-validated
      walk gets a consistent view by retrying; here a reader never retries,
      and what it is spared is exactly new-then-old --- not the straddle.
\item Strictly more than independent registers, strictly less than a
      snapshot --- and the middle ground is usable, not merely tolerable.
\item [opt] Split edits: Triplett's rule, commit in reverse of the reader's
      read order (backup slide).
\item If asked: two slots reached independently (a cached pointer, sibling
      subtrees) have no chain and owe an acquire; it bites under
      -DURCU\_DEREFERENCE\_USE\_VOLATILE on weakly-ordered hardware.
\item If asked, kernel side: arm64 with CONFIG\_LTO already promotes
      READ\_ONCE() to acquire, because LTO may turn an address dependency
      into a control one.
\end{itemize}}
\end{frame}

\begin{frame}{Between RCU and STM}
\begingroup\renewcommand{\arraystretch}{1.3}\arrayrulecolor{muted!60}
\begin{tabular}{@{}>{\raggedright}p{31mm}>{\raggedright}p{26mm}>{\columncolor{accent!9}\raggedright}p{40mm}>{\raggedright\arraybackslash}p{30mm}@{}}
\toprule
& \textbf{RCU} & \textbf{\color{accent}Pseudo-transaction} & \textbf{Opaque STM} \\
\midrule
Atomic to readers & one slot & N slots, one commit & read and write sets \\
The reader pays & nothing & one predicted branch; no added load, no
  per-element state & a barrier on each transactional read \\
The reader gets & one slot at a time & never new-then-old
  & isolation, opacity \\
\bottomrule
\end{tabular}\par\endgroup
\vspace{4mm}
\lead{Isolation is a read-side cost once updates mutate in place.
\textbf{RCU pseudo-transactions decline to charge readers for it.}}
\vspace{2mm}
\aside{Honest asterisk: RLU gives readers a snapshot, existence a per-element
membership test --- strictly more, and charged to readers
(slide~\ref{fr:comparison}).}
\note{%
\begin{itemize}
\item ``Once updates mutate in place'' is load-bearing: copy-on-write gives
      readers a free snapshot, which is why a reader gets a coherent node for
      free. Copying just does not scale to the update.
\item We did not fail to provide isolation; we declined to charge for it.
      A reader who needs a snapshot needs a different facility.
\item The column is an opaque STM (TL2 is one), not every STM. If asked: a
      read can opt out of the barrier, and is then not isolated. Reads outside
      any transaction under weak atomicity; Intel's \code{[[waiver]]} and
      \code{transaction\_pure}; STAMP's hand instrumentation of ``important''
      loads only (Kestor et al., TRANSACT 2011).
\item Closest on that side: Howard and Walpole (HotPar 2011) run RCU readers
      outside a modified SwissTM. Their readers ``can see partially completed
      commits'': the commit replays its writes in program order, grace
      periods included. Here the reader pays one branch, and the commit is
      atomic to it.
\end{itemize}}
\end{frame}

% ===========================================================================
% D, continued. The API: the structures built on it
% ===========================================================================
\talksection{The API}

\begin{frame}[fragile]{The list: the reverse walk comes back}
\begin{codeblock}
rcu_read_lock();
urcu_txn_sw_list_for_each_entry_reverse_rcu(pos, &head, node)
    ...;
rcu_read_unlock();

urcu_txn_sw_list_del_rcu(&e->node);    /* writer, under the embedder's locks */
call_rcu(&e->rcu, free_e);
\end{codeblock}
\vspace{4mm}
\lead{At every resolution, \code{a->next == b} $\Rightarrow$
\code{b->prev == a}.}
\vspace{3mm}
\begin{itemize}\small
\item A removed node keeps its own \code{next} and \code{prev} --- a ghost ---
      so a reader standing on it escapes either way.
\item Coherence holds \emph{at} each resolution. Two resolutions that must
      agree are asking for a snapshot, which is not on offer.
\end{itemize}
\note{%
\begin{itemize}
\item Forward and backward chains describe the same list at every instant:
      the reverse walk rculist declines becomes one a reader may take.
\item A prev hop taken from the node the previous load returned is
      dependency-chained. Comparing a prev and a next read from nodes the
      reader did not chain between needs its own acquire.
\end{itemize}}
\end{frame}

\begin{frame}[fragile]{What a list edit records: two edges, one flip}
\begin{codeblocksmall}
static inline int
urcu_txn_sw_list_add_after_prepare(struct urcu_txn_sw_txn *txn,
        struct urcu_txn_sw_list_node *newp,
        struct urcu_txn_sw_list_node *pos)
{
    struct urcu_txn_sw_list_node *next =
        urcu_txn_sw_list_pending(txn, &pos->next);
    newp->prev = pos;                   /* build: unreachable */
    newp->next = next;
    urcu_txn_sw_record_chain(txn, (void **) &pos->next, next, newp,
                             URCU_TXN_SW_LIST_PROXY_TAG);
    urcu_txn_sw_record_chain(txn, (void **) &next->prev, pos, newp,
                             URCU_TXN_SW_LIST_PROXY_TAG);
    return 0;
}
\end{codeblocksmall}
\vspace{1mm}
{\small The \code{\_rcu} wrappers: init + declare\_disjoint + reserve(2) +
\code{\_prepare} + commit.\par}
\aside{Abridged from \code{rcu-txn-sw-list.h}: \code{(void)} casts and
comments dropped, lines re-wrapped.}
\note{%
\begin{itemize}
\item The neighbour is read through the pending view: an earlier edit in the
      same transaction may already have moved pos->next.
\item Build is private: newp's own links are set before it is reachable.
\end{itemize}}
\end{frame}

\begin{frame}[fragile]{One commit, two structures}
\begin{codeblock}
urcu_txn_sw_init(&txn);
urcu_txn_sw_reserve(&txn, 3);          /* 1 hlist edge + 2 list edges */

/* findable: */
urcu_txn_sw_hlist_add_head_prepare(&txn, &e->hnode, bucket);
/* ordered: */
urcu_txn_sw_list_add_after_prepare(&txn, &e->lru, &lru.node);

ret = urcu_txn_sw_commit(&txn);        /* one flip */
\end{codeblock}
\vspace{3mm}
\lead{$E$ becomes findable \textbf{and} ordered at the same instant.
Hiding it has the same shape.}
\vspace{2mm}
\aside{The abstract's composability claim: one object in several
transaction-aware structures, made visible or hidden atomically. Slide
\ref{fr:skew}'s composition problem has no middle state left.}
\label{fr:twostructs}
\note{%
\begin{itemize}
\item Reserve before composing: the hlist \_prepare asserts (debug build) that
      the bracket stays rollbackable.
\item Two commits would be two linearization points again: the facility makes
      the composition available, not automatic.
\end{itemize}}
\end{frame}

\begin{frame}{Read-your-own-writes}
\begin{columns}[T, onlytextwidth]
\column{0.47\textwidth}
  \vspace{2mm}
  \centerfig{fig-ryw}
  \vspace{5mm}
  \begingroup\small\arrayrulecolor{muted!60}\renewcommand{\arraystretch}{1.2}
  \begin{tabular}{@{}lll@{}}
  \toprule
  \textbf{slot} & \textbf{old} (committed) & \textbf{new} \\
  \midrule
  \code{P.next}  & E1 & {\color{muted}E2 →} \textbf{\color{accent}N}
                        \ {\scriptsize\color{muted}fused} \\
  \code{E2.prev} & E1 & P \\
  \code{N.prev}  & E2 & P \\
  \bottomrule
  \end{tabular}\par\endgroup
\column{0.5\textwidth}
  \small
  The second delete must see the first one's \emph{pending} \code{P.next}.
  \begin{itemize}
  \item \code{urcu\_txn\_sw\_load()} returns this transaction's pending value.
  \item \code{record\_chain()} fuses onto the record: keeps the committed old,
        advances the new.
  \item Find: a linear scan of a private array; a Bloom filter armed lazily
        for long write sets.
  \end{itemize}
  \vspace{1mm}
  Without it, the second delete records \code{E1.next} on a node about to
  vanish.
\end{columns}
\vspace{4mm}
\aside{Three obligations: search through the transaction's load, or before the
first record; every anchor still live; \code{declare\_disjoint()} is a promise
(\code{-DURCU\_TXN\_SW\_DEBUG\_DISJOINT} traps a broken one). Read-your-own-writes
itself is standard in deferred-update STM.}
\note{%
\begin{itemize}
\item Nothing here is a race: one writer, editing in program order, would
      tear its own commit.
\item The load records nothing: no read set, validated against nothing.
      Under exclusion there is nothing to validate against.
\item A reader is not in the scheme at all: the writer chooses per load what
      an STM applies wholesale, and the reader is what the choice buys.
\end{itemize}}
\end{frame}

\begin{frame}{The hlist: one transacted edge}
\lead{Readers walk an hlist forward only, so an operation has \textbf{one}
reader-visible edge.}
\vspace{4mm}
\begin{itemize}
\item \code{pprev} is writer-only and plain-stored: no reader loads it, and no
      peer races it.
\item A lone insert or delete costs about an \code{rcu\_assign\_pointer()}: no
      proxy, no group, no \code{call\_rcu()}.
\item It still composes, through \code{\_prepare}: the single edge folds into a
      larger flip (slide~\ref{fr:twostructs}).
\item Single-pointer head: a bucket costs 8 bytes, and no operation
      special-cases the head.
\end{itemize}
\note{%
\begin{itemize}
\item [opt] Cut first if short on time.
\item pprev is written eagerly, in the prepare step. Deferring it and reading
      it raw would be the hlist form of the adjacent-delete trap.
\end{itemize}}
\end{frame}

% ===========================================================================
% E. Cost
% ===========================================================================
\talksection{Cost}

\begin{frame}{Reader: zero loads, one branch}
% One list against the other, per node (P1 sec. 7: the same eight
% instructions, plus one test and one branch, and no load).
\begingroup\small
\begin{tabular}{@{}l@{\hspace{3mm}}l@{\hspace{3mm}}r@{ }l@{\hspace{5mm}}l@{}}
Per node: & \code{rculist}     & 8  & instructions: 2 branches, 2 loads & \\
          & \code{txn\_sw\_list} & 10 & instructions: 3 branches, 2 loads
          & \textbf{+1 test, +1 branch, 0 loads} \\
\end{tabular}\par
\endgroup
\vspace{1mm}
{\small\color{muted}Read throughput relative to \code{rculist} (dashed = 1.0),
\textbf{higher is better}; 192 readers, no writer.\par}
\centerfig{fig-readcost}
\vspace{2mm}
\begin{itemize}\small
\item \textbf{Scattered nodes}, as allocated over time: the branch is
      \textbf{free}, 1.002--1.003 of \code{rculist}.
\item \textbf{Walk order, L1-resident}: it shows, 10.7\% (the control: the same).
\end{itemize}
% The setup of every cost slide (P1 sec. 7), stated once, on the first;
% the rest of it is in the note.
\aside{Setup: 2× AMD EPYC 9654 (192 cores), one thread per core; liburcu QSBR;
10\,000 nodes; median of five 3\,s runs.}
\note{%
\begin{itemize}
\item The ratio is of throughput (node visits/s, no writer: P1's ``read
      ceiling''), each arm against rculist in the same run. 0.893 means
      10.7\% slower. Every curve is flat in the reader count.
\item Scattered: matches the +tag test control within 0.05\%. L3- and
      DRAM-sized scattered lists: no arm distinguishable from rculist within
      run-to-run noise (13\% at L3, 5\% at DRAM).
\item Always quote with the layout. The 10.7\% is a property of a list the
      prefetcher keeps in L1, walked in allocation order.
\item Read ceiling vs rculist, no writer, 10\,000 nodes walked forward twice
      by every arm; median of five 3 s runs after 4 s warm-up. Every curve is
      flat in the reader count; scaling 97.1--98.4\% of ideal.
\item Controls: +tag test has the facility's exact instruction stream, and
      reads the same in both layouts (10.4--11.2\% against 10.6--10.8\% in
      walk order). An earlier take had the two 2\% apart there: code layout,
      not mechanism. Few-\% gaps between the load controls are code-level
      noise.
\item Setup: 2× EPYC 9654, 192 cores, liburcu QSBR, per-CPU call\_rcu
      workers, jemalloc per-CPU arenas.
\end{itemize}}
\end{frame}

\begin{frame}{Against schemes that give readers more}
\begingroup\small\renewcommand{\arraystretch}{1.4}\arrayrulecolor{muted!60}
\begin{tabular}{@{}llcclc@{}}
\toprule
& \textbf{Per-element} & \textbf{Extra} & \textbf{Extra}
  & \textbf{Reader} & \textbf{Reader pays vs txn} \\[-1.5mm]
& \textbf{state} & \textbf{loads} & \textbf{branches}
  & \textbf{guarantee} & {\scriptsize\color{muted}scattered / walk order} \\
\midrule
Existence & \code{eh\_egi} field & 1 & 1 & per-element & 2--3\% / 15--17\% \\
RLU & object header & 1 & 1 & snapshot & 1--3\% / 23--44\% \\
MV-RLU & header + versions & $\ge 2$ & $\ge 3$ & snapshot & 1.43× / 3.01× \\
\midrule
\rowcolor{accent!10}
\textbf{This work} & \textbf{none} & \textbf{0} & 1 & never new-then-old & --- \\
\bottomrule
\end{tabular}\par\endgroup
\vspace{7mm}
\lead{They charge readers more \textbf{and} deliver more: a price, not a
defect.}
\vspace{3mm}
\aside{Existence in its best layout (a 32-byte element); MV-RLU in its
global-clock configuration, ratios at 192 readers. Write throughput is not
ranked across this boundary.}
\label{fr:comparison}
\note{%
\begin{itemize}
\item Every design branches, this one included: the branch is not where the
      comparison is decided. The others must first load what they test; ours
      tests a value the reader already holds. The 1 is not bold on purpose.
\item RLU's pure reader never touches the header: its extra load is a
      per-thread flag. Do not claim ``RLU loads the header on every read''.
\item RLU walk order: 23--27\% up to 32 readers, then bends above 64 (85\% of
      ideal at 192), unexplained.
\item Scattered figures are from two readers up. At one reader existence and
      RLU read low in four runs of five this take (7\% and 9\% gaps there);
      not explained, not quoted.
\item The unconditional cost of the others is space: a header or field on
      every element, for as long as the structure lives.
\end{itemize}}
\end{frame}

\begin{frame}{Writer: 1.3--1.8×, and it is not the flip}
\begin{columns}[T, onlytextwidth]
\column{0.27\textwidth}
  {\fontsize{25}{28}\selectfont\bfseries\color{accent}1.3--1.8×\par}
  \vspace{1mm}
  {\footnotesize\color{muted}one writer, readers running\par}
\column{0.7\textwidth}
  \vspace{1mm}
  \begin{itemize}\small
  \item Scattered nodes: 1.55--1.58× up to 4 readers, 1.6--1.8× beyond.
  \item Walk order: 1.34--1.39× up to 4 readers, 1.26--1.52× beyond.
  \item No reader: no ratio. Plain RCU outruns its own reclaim there.
  \end{itemize}
\end{columns}
\vspace{1mm}
{\small\bfseries Where the cycles go\enspace
 {\mdseries\color{muted}one writer, one reader, no mutex: the writer's thread,
 per update}\par}
\vspace{1mm}
\centerfig{fig-profile}
\vspace{1mm}
\aside{Intrinsic: a fresh txn descriptor per update, kept for a grace period.
The commit: 19 cycles of 243.}
\aside{Not intrinsic: this slab. Under rseq: no lock, no atomics. A kernel
port could use the allocator and \code{kfree\_rcu()}.}
\note{\footnotesize	% at the notes' own size the last lines fall off the page
\begin{itemize}
\item A parked group block must outlive a grace period: a reader that loaded
      a proxy may still be resolving through it. Plain RCU defers only the
      node. A single-edge commit needs no proxy and frees in place.
\item At engine 2793224e (3 Oct). Earlier copies said about 2×, at 18809ea8:
      no batched retirement, an out-of-line append, a 416-byte descriptor
      (now 224), a redundant full barrier in the slab's pop and push.
\item The slab's 95 cycles: over half on four locked instructions (pop lock
      and unlock, a cmpxchg per pop and per push), the rest on a block that
      comes back cache-cold a grace period later. The four: a freelist any CPU
      may push to, not the mechanism; the rseq slab has none, 1.5× scattered,
      1.3× walk order up to 4 readers, same from 8 up. Not profiled. Kernel
      (not P1's, not measured): the allocator and kfree\_rcu(), no slab.
\item Batching is worth 9--23\% to the writer. With one call\_rcu per
      descriptor (URCU\_TXN\_SLAB\_NO\_BATCH) the ratio is 1.9--2.0×
      scattered, 1.4--1.7× walk order.
\item No reader: plain RCU's writer defers nodes faster than the call\_rcu
      worker on its own hardware thread frees them: the worker never sleeps
      and the process grows by more than 100 MB/s, a rate it cannot sustain.
      With 1 to 8 readers both lists hold memory flat. Earlier slides quoted
      2.5--2.6× there: do not.
\item Both collapse from 4 to 8 readers (longer grace periods): scattered,
      RCU 14.3 → 4.7 Mops/s, txn 9.3 → 2.6.
\item Node allocation reads lower for the transacted writer (21 against 50).
      Not explained; do not claim it.
\item We report with readers: a structure with no readers has no use for
      reader-visible coherence.
\end{itemize}}
\end{frame}

\begin{frame}{Writers scale with fine-grained locks}
\begingroup\small Writers on disjoint slots: plan unlocked → lock in ascending address
order → validate → commit.\par
\endgroup
\vspace{1mm}
% P1's fig:writerscale, restyled for projection: direct labels at the line ends
% instead of a legend, the endpoints the talk quotes printed there, and the
% handoff's CVD-checked series colours (HANDOFF.md section 4), each also told
% apart by its dash pattern. Data read in place from P1.
%
% The bars are P1's min-max over five runs and stay: the spread at 128-160
% writers is real (memory placement across sockets), and hiding it would make
% those medians look tighter than they are.
\pgfplotstableread[col sep=comma]{\ponedata/fig-writerscale.csv}\wstab
\lastof{\wstab}{bit}{\wsBit}       \lastof{\wstab}{stripe}{\wsStripe}
\lastof{\wstab}{rcu}{\wsRcu}       \lastof{\wstab}{mutex}{\wsMutex}
\def\wsnum#1{\pgfmathprintnumber[fixed, precision=0]{#1}}
\begin{tikzpicture}
\begin{axis}[
  scale only axis, width=9.4cm, height=3.0cm,
  xlabel={Writers (disjoint slots, no readers)},
  ylabel={Millions of updates/s},
  xmin=0, xmax=196, ymin=0, ymax=420,
  xtick={1,32,64,96,128,160,192}, ytick={0,100,200,300,400},
  tick label style={font=\footnotesize, text=muted},
  label style={font=\footnotesize, text=muted},
  axis lines=left, axis line style={muted!60}, tick style={muted!60},
  ymajorgrids, grid style={muted!18},
  clip=false,
  every axis plot/.append style={line width=1.3pt, mark size=1.2pt},
  error bars/error bar style={line width=0.5pt, solid},
  error bars/error mark options={rotate=90, mark size=1.2pt, line width=0.5pt},
]
\draw[muted!70, densely dotted] (axis cs:96,0) -- (axis cs:96,420)
  node[anchor=north west, font=\scriptsize, text=muted] {second socket →};

\addplot[color=onelock, dotted, mark=triangle*,
         error bars/.cd, y dir=both, y explicit]
  table[x=writers, y=mutex, y error plus=mutex_ep, y error minus=mutex_em,
        col sep=comma] {\ponedata/fig-writerscale.csv};
\addplot[color=old, solid, mark=square*,
         error bars/.cd, y dir=both, y explicit]
  table[x=writers, y=rcu, y error plus=rcu_ep, y error minus=rcu_em,
        col sep=comma] {\ponedata/fig-writerscale.csv};
\addplot[color=stripes, dashed, mark=o,
         error bars/.cd, y dir=both, y explicit]
  table[x=writers, y=stripe, y error plus=stripe_ep, y error minus=stripe_em,
        col sep=comma] {\ponedata/fig-writerscale.csv};
\addplot[color=accent, solid, mark=*, line width=1.8pt,
         error bars/.cd, y dir=both, y explicit]
  table[x=writers, y=bit, y error plus=bit_ep, y error minus=bit_em,
        col sep=comma] {\ponedata/fig-writerscale.csv};

% Direct labels, each printing its series' median at 192 writers (the CSV's
% last row). rculist and the bit lock end 2 apart (379 vs 377), so their labels
% are pulled apart vertically.
\node[anchor=west, font=\footnotesize, text=old]
  at (axis cs:197,408) {\code{rculist}, stripes: \wsnum{\wsRcu}};
\node[anchor=west, font=\footnotesize\bfseries, text=accent]
  at (axis cs:197,362) {\code{txn\_sw}, lock bit: \wsnum{\wsBit}};
\node[anchor=west, font=\footnotesize, text=stripes]
  at (axis cs:197,302) {\code{txn\_sw}, stripes: \wsnum{\wsStripe}};
\node[anchor=west, font=\footnotesize, text=onelock]
  at (axis cs:197,14) {one lock: ≈\wsnum{\wsMutex}};
\end{axis}
\end{tikzpicture}

\vspace{1mm}
{\small Lock bit in each node: \textbf{10.3 → 386}, 38× from 1 to 192 writers,
and 339× the list behind one lock. rculist cannot use an in-node lock: its
\code{prev} cannot be accessed coherently, so finding the predecessor takes a
forward walk.\par}
\label{fr:writerscale}
\note{%
\begin{itemize}
\item rculist leads 1.57× alone, 1.41× at 2 writers and 1.03--1.28× from 4
      to 96; level from 128 on (0.94--1.04): 386 vs 403 at 192. Under the same stripes:
      1.6× alone, 0.89--1.35 from 16 writers up.
\item To use an in-node lock, rculist would first have to find the
      predecessor by a forward walk; this facility's writer reads it from the
      victim's own prev, as a reader may.
\item Caveats to say: each writer owns its slots, so no retries; no lookup
      cost in the figure; memory placement matters on this box --- up to 96
      writers memory on socket 0, beyond interleaved; 128--160 bistable
      (bars, up to 21\%), 191 → 192 jump (5--19\%) accepted.
\end{itemize}}
\end{frame}

% ===========================================================================
% F. Kernel and close
% ===========================================================================
\talksection{Kernel and close}

\begin{frame}{Dentry cache: bucket locks around SW commits}
\begin{columns}[T, onlytextwidth]
\column{0.34\textwidth}
  \begingroup\small
  A \textbf{userspace model}, two engines.\par
  \vspace{1mm}
  \textbf{\color{old}The kernel's scheme.}
  \code{d\_seq} hand over hand, \code{rename\_lock}, \code{i\_rwsem}.\par
  \vspace{1mm}
  \textbf{\color{accent}Bucket lock + SW txn.}
  Bit locks on the bucket and child-list heads. A rename is \textbf{one
  commit} across both indexes. No \code{d\_seq}, no \code{rename\_lock}, no
  lock in \code{readdir}.\par
  \endgroup
  \vspace{1.5mm}
  \aside{Reader rows: both engines at the same paced rename load.
  A userspace model is not kernel evidence: no reference counts, and no
  \code{i\_rwsem}, which filesystems rely on.}
\column{0.65\textwidth}
  \centerfig{fig-dcache}
\end{columns}
\label{fr:dcache}
\note{%
\begin{itemize}
\item The shape is slide \ref{fr:twostructs}'s: a rename leaves the old hash bucket and child
      list and enters the new ones in one flip, under fine-grained locks
      (slide \ref{fr:exclusion}).
\item If asked which API: the model commits through \code{rcu-txn.h}'s
      single-writer path, \code{urcu\_txn\_store\_sw()} and
      \code{urcu\_txn\_commit\_sw()}: park by plain store, one release store
      of a status word, settle; no CAS, no contention abort. It is not
      \code{rcu-txn-sw.h}'s \code{urcu\_txn\_sw\_commit()}, and its readers
      resolve through the record's descriptor, not a proxy's group.
\item 11 of the benchmark summary's 22 rows. The rows cut are all on par or
      faster; both of its slower rows are here. Sweep
      \texttt{\dcacheSweep}, benchmark tree \texttt{\dcacheBench}; a dot is
      one measured point.
\item A reader row counts a point only where both engines' writers kept the
      pace: the reverse walk's \dcval{dpath}{hi}× is at 32 readers, and from
      48 on the baseline's renamers starve.
\item Why it loses, and what the writers' lead includes: backup slide
      \ref{fr:dcachemodel}.
\item The whole-path consistency the model's lookups get is the embedder's
      re-check of a deletion mark, not the facility's: slide \ref{fr:provides} stands.
\item Say ``model'' and ``userspace'' each time: that \code{rename\_lock} and
      \code{d\_seq} can go is the claim most likely to be challenged.
\end{itemize}}
\end{frame}

\begin{frame}{What a kernel port needs}
\placeholder{Port status --- Mathieu to confirm and fill.}
\vspace{1.5mm}
\begin{columns}[T, onlytextwidth]
\column{0.48\textwidth}
  {\bfseries Needs\par}
  \vspace{1mm}
  \begin{itemize}\small
  \item A spare low bit in the slot: list pointers are aligned.
  \item A commit fails only on OOM, before anything is parked. Reserve
        before the first irreversible step and it cannot fail --- under a
        spinlock, say.
  \item One \code{call\_rcu()} per multi-edge commit (batchable).
  \item The exclusion validator as a debug option.
  \end{itemize}
\column{0.48\textwidth}
  {\bfseries Where it would land\par}
  \vspace{1mm}
  \begin{itemize}\small
  \item A coherent reverse walk for rculist: the iterator it declines
        today.
  \item Cross-structure publish: hash + LRU.
  \item Pinned objects whose back edges must move in place.
  \end{itemize}
  \vspace{1mm}
  {\bfseries For the dentry cache, also\par}
  \vspace{0.5mm}
  \begin{itemize}\small
  \item Reference counts, or hazard pointers?
  \item \code{i\_rwsem}: filesystems rely on it.
  \end{itemize}
\end{columns}
\label{fr:kport}
\note{%
\begin{itemize}
\item The dentry-cache block is about the model of slide \ref{fr:dcache},
      not about the facility. No count: every
      walk stays inside one RCU read-side section and a handle is valid only
      while the caller keeps the object from being unlinked or evicted. The
      kernel leaves the RCU walk through \code{try\_to\_unlazy()}, which takes
      \code{d\_lockref} before anything blocks, and an open file, a working
      directory and a cached child each hold one. Whether a count is the
      right tool for that, or whether this is a use for hazard pointers, is
      not evaluated.
\item \code{i\_rwsem} is the VFS's directory lock, not the dentry cache's:
      \code{Documentation/filesystems/locking.rst} has \code{lookup} under it
      shared and \code{create}, \code{link}, \code{unlink}, \code{mkdir},
      \code{rmdir}, \code{rename} exclusive, and the VFS takes it before it
      calls the filesystem. The model shows the dentry cache can stop needing
      it; it shows nothing about a filesystem that does.
\item P1 is a userspace paper and says nothing about a kernel port; every item
      here maps a P1 statement onto the kernel. The failure and reserve line
      is P1 sec.\ 4.5 (``before the first irreversible side effect''); ``under
      a spinlock'' is the kernel reading of it, not P1's.
\item Do not name list\_bidir's caller as a customer: P1 sec.\ 2.4 says it
      takes one hop and never needs coherence.
\end{itemize}}
\end{frame}

\begin{frame}{What it is not}
\begin{itemize}
\item \textbf{Not STM.} No snapshot, no opacity: a reader is not a transaction
      of any kind.
\item \textbf{Exclusion is the embedder's}, per slot, with its own locks.
      Two writers on one slot corrupt it.
\item \textbf{The tag contract}: no legitimate slot value may carry the tag
      pattern. The contracts underwrite memory safety, not just atomicity.
\item \textbf{Commit width grows with edges}: a tall tower commits many
      records where existence flips one group.
\item \textbf{The writer pays 1.3--1.8×.}
\end{itemize}
\note{%
\begin{itemize}
\item These are properties of the model or contracts on the embedder, not
      defects: the weakening is targeted, not residual.
\end{itemize}}
\end{frame}

\begin{frame}{Summary}
\begin{itemize}
\item RCU gives a writer one atomic store. A pseudo-transaction widens it to
      \textbf{N slots, still one store}: the flip-latch.
\item Readers: \textbf{no added load, no per-element state}; one predicted
      branch, free on scattered nodes.
\item Writers: 1.3--1.8×, mostly the txn descriptor's slab, not the flip. They scale
      with the embedder's fine-grained locks.
\item Never new-then-old, and no more: \textbf{not STM}.
\end{itemize}
\vspace{1mm}
\begin{callout}
\textbf{Paper}\enspace \ph{arXiv link --- TBD}\\[1mm]
\textbf{Code}\enspace \code{github.com/compudj/userspace-rcu-dev} @
\code{2793224e}, in \code{include/urcu/}:\\
\code{rcu-txn-sw.h}, \code{rcu-txn-sw-list.h}, \code{rcu-txn-sw-hlist.h},
\code{rcu-txn-status.h}
\end{callout}
\vspace{1mm}
{\Large\bfseries Questions?\par}
\end{frame}

% ===========================================================================
% Backup
% ===========================================================================
\appendix
\talksection{Backup}

\begin{frame}{Anticipated questions}
\footnotesize
\begin{columns}[T, onlytextwidth]
\column{0.48\textwidth}
  \textbf{Why not a seqcount?}\\
  A read-side retry on every traversal, detecting after the fact. Here a
  reader never retries --- and gets no snapshot either.
  \vspace{0.8mm}

  \textbf{Why not copy?}\\
  The cost scales with the region, and a pinned node cannot be copied.
  \vspace{0.8mm}

  \textbf{Isn't this existence?}\\
  Existence adds a per-element field (+1 load, +1 branch) and gives a
  stronger per-element guarantee. Pseudo-transactions add no per-element
  state; existence's readers pay 3\% of read throughput on scattered nodes,
  15\% in walk order.
  \vspace{0.8mm}

  \textbf{ABA?}\\
  Resolution goes through the descriptor, not the slot value: slot-value
  A-B-A is inert.
\column{0.48\textwidth}
  \textbf{HTM?}\\
  Commodity HTM may abort, so it needs a fallback that meets the same problem.
  s390 constrained transactions: at most 32 instructions in 256 bytes.
  \vspace{0.8mm}


  \textbf{Does the tag test need a barrier?}\\
  No. The proxy's immutable fields are reached through the dependency; only
  the selector is loaded with acquire, and only when parked.
  \vspace{0.8mm}

  \textbf{What if two writers race?}\\
  Silent corruption, by contract. \code{-DURCU\_TXN\_SW\_EXCL\_VALIDATE}
  catches it.
  \vspace{0.8mm}

  \textbf{Isn't this MCAS?}\\
  The descriptor-and-status shape is: Harris et al.\ (2002). Nearest,
  Guerraoui et al.\ (2020): $k{+}1$ CAS, unlock deferred to reclamation.
  Here every step is a plain store, and settle is eager.
\end{columns}
\note{%
\begin{itemize}
\item MCAS, in full. Guerraoui, Kogan, Marathe, Zablotchi, ``Efficient
      Multi-word Compare and Swap'' (DISC 2020, arXiv 2008.02527): lock-free,
      helping, plain-CAS install with no RDCSS, $k+1$ CAS, unlock deferred to
      reclamation by epochs they call similar to RCU, readers that do not
      write unless they meet an in-flight operation, and a proof that a
      lock-free disjoint-access-parallel $k$-CAS must CAS at least $k$
      locations.
\item Against this talk's engine: zero CAS, bought by exclusion, which is
      outside that bound's hypotheses; and eager settle. P1 sec.\ 11.1 and
      9.4.
\item The abstract says ``novel''. Not the descriptor mechanism. What P1
      sec.\ 11.3 says it has not found: a marker that is transient, that RCU
      readers resolve through without waiting or writing, and that adds no
      per-element state, with the per-record tag that follows.
\item The multi-writer engine is not in this talk. If asked: the same
      descriptor with a CAS install and no helping. It is the engine their
      paper is closest to: both install by plain CAS with no RDCSS. It differs
      by settling at once and deciding with a store, and it is not lock-free;
      theirs is.
\end{itemize}}
\end{frame}

\begin{frame}{Split edits: publish in reverse read order}
\lead{When one edit must span several commits, commit in the \textbf{reverse}
of the reader's read order.}
\vspace{4mm}
\begin{itemize}
\item Reader resolves $s_1, \ldots, s_k$ in that order; the writer publishes
      $s_k$ first, $s_1$ last. Once a reader sees one slot new, every later
      slot is new too: a prefix of old, a suffix of new.
\item Example: re-parent $E$ from $P$ to $Q$. Commit \code{E->parent} first,
      then the child links. A reader that descends into $E$ under $Q$ then
      finds \code{E->parent == Q}.
\item Needs one fixed read order and distinct slots. A list read in both
      directions has none to reverse: it needs the atomic commit.
\end{itemize}
\vspace{3mm}
\aside{First stated, to our knowledge, by Triplett. An embedder's discipline:
the facility cannot know a structure's traversal order.}
\end{frame}

\begin{frame}{One writer, with readers}
\centerfig{fig-writecost}
\vspace{4mm}
{\small Both writers collapse between four readers and eight. Scattered: the
ratio is 1.55--1.58 up to four readers and 1.61--1.80 beyond. Walk order:
1.34--1.39 up to four readers, 1.30 at eight, 1.26--1.52 up to 191.\par}
\note{Uncontended writer mutex on both; both reclaim through call\_rcu(), same
RCU flavor, so the separation isolates the pseudo-transaction layer. Median of
five 3 s runs after 4 s warm-up.}
\end{frame}

\begin{frame}{Dentry cache: the model, and where it loses}
\footnotesize
\begin{columns}[T, onlytextwidth]
\column{0.46\textwidth}
  \textbf{The baseline is the kernel's}
  \begin{itemize}
  \item Lookup: the RCU walk. \code{d\_seq} hand over hand;
        \code{rename\_lock} sampled at \code{path\_init()}, consulted on a
        miss.
  \item Reverse walk: a port of \code{dentry\_path\_raw()}.
  \item Locks: the kernel's \code{rw\_semaphore} as \code{i\_rwsem}, bit
        locks that spin as \code{bit\_spin\_lock()} does, a queued spinlock
        on the LRU shards.
  \end{itemize}
  \vspace{0.5mm}
  \textbf{The comparison}
  \begin{itemize}
  \item Readers are compared under paced writers, at the points both
        engines' writers sustain.
  \end{itemize}
  \vspace{0.5mm}
  \textbf{Not modelled}
  \begin{itemize}
  \item Reference counts (\code{lockref}): walks stay in RCU. Mounts, hard
        links, security hooks.
  \end{itemize}
\column{0.52\textwidth}
  \textbf{How the model commits}
  \begin{itemize}
  \item \code{rcu-txn.h}'s single-writer commit: park, one flip, settle.
        No CAS, no contention abort.
  \item Not \code{rcu-txn-sw.h}: readers resolve through the commit's
        descriptor, not a proxy's group.
  \end{itemize}
  \vspace{0.5mm}
  \textbf{Where it loses}
  \begin{itemize}
  \item Hits on objects being renamed: \dcrange{hit}. A rename publishes a
        shell; until a \code{call\_rcu} callback folds it back, about 5 ms, a
        hit reads three cache lines instead of one.
  \item At rest: \dcrange{hit-rest}, per-hop pointer decoding.
  \end{itemize}
  \vspace{0.5mm}
  \textbf{What the writers' lead includes}
  \begin{itemize}
  \item The baseline takes \code{i\_rwsem}, which filesystems rely on, and
        \code{s\_vfs\_rename\_mutex}; the model neither. A port gains only
        what the VFS stops taking.
  \end{itemize}
\end{columns}
\label{fr:dcachemodel}
\note{%
\begin{itemize}
\item Sweep \texttt{\dcacheSweep} (2026-10-03), benchmark tree
      \texttt{\dcacheBench}: \code{experiments/dcache/README.md}, Results.
\item On the far socket the baseline's readdir and reverse walk lose
      30--40\%, the model 2--5\%: its readers write no shared line.
\item Every run ends on a namespace-conservation check; a run that fails it
      is not plotted. 200 ms of the real workload before each timed window.
\item A point counts for a reader row only if both engines' writers kept the
      offered pace. The namespace sits on the benchmark's own NUMA node, the
      baseline's favourable case. The model needs neither lock because its \code{readdir}
      takes none: no writer has a reader to exclude. That is the dentry
      cache's need for \code{i\_rwsem}, not the filesystem's (slide
      \ref{fr:kport}).
\item No count in either engine, so neither pays the reference a kernel walk
      takes on its last component. The eviction model stands in for it: a
      directory with cached children is never evicted.
\item \code{urcu\_txn\_store\_sw()} records, \code{urcu\_txn\_commit\_sw()}
      commits; a reader resolves with \code{urcu\_txn\_resolve\_record()}:
      the descriptor's status word picks old or new.
\end{itemize}}
\end{frame}

\end{document}
